\section{결론}

현재 결과를 기준으로 FOKVQ 논문의 가장 안전하고도 의미 있는 개정 방향은 분명하다. 이 논문은 더 이상 ``PCA 기반 양자화 방법이 baseline을 이긴다''를 중심 주장으로 삼기보다, \textbf{회전 기반 KV-cache 압축을 Lie 군 위의 좌표계 선택 문제로 해석하고, 그 안에서 PCA가 왜 합리적인 정적 회전인지 수학적으로 정당화하며, 그 해석이 empirical surrogate에서 실제로 어디까지 지지되는지 보이는 논문}이 되어야 한다.

이 방향은 이미 확보된 사실과 맞는다. 비등방성과 facet 구조는 실재한다. 정적 회전 선택 문제에서는 PCA가 강한 기준선이다. query-dependent 변화는 존재하지만, 축 자체보다 스케줄 문제를 더 강하게 시사한다. 반면 현재 구현의 low-bit PPL 우위는 아직 충분히 입증되지 않았다. 그러므로 가장 정직한 결론은 ``Lie 군 해석은 지지되지만, stronger quantizer까지 포함한 완결된 시스템 우위는 후속 검증이 필요하다''이다.

이 재구성은 논문을 약하게 만드는 것이 아니다. 오히려 현재 데이터가 실제로 지지하는 수준에 맞춰 claim hierarchy를 다시 세우는 일이다. 그 위에서 추가 실험이 성공하면 practical quantizer paper로 확장할 수 있고, 그렇지 않더라도 회전 선택과 정보 집중 구조에 대한 해석 논문으로서 충분한 의미를 유지할 수 있다.
